% from ega0-18-fr.tex
\oldpage[0\textsubscript{IV-1}]{63}

on en déduit canoniquement des homomorphismes
\begin{equation}
\label{0.18.3.7.2-fr}
\tag{18.3.7.2}
  \operatorname{Exan}_{A/\Lambda}(B,L)\longrightarrow
  \operatorname{Exan}_{A'/\Lambda}(B,L)\longrightarrow
  \operatorname{Exan}_{A'/\Lambda}(B',L)
\end{equation}
qui proviennent de la commutativité du diagramme
\[
  \xymatrix@C=2.8pc@R=2.3pc{
    \operatorname{Exan}_A(B,L)\ar[r]\ar[d]
      & \operatorname{Exan}_{A'}(B,L)\ar[r]\ar[d]
      & \operatorname{Exan}_{A'}(B',L)\ar[d]\\
    \operatorname{Exan}_{\Lambda}(B,L)\ar[r]
      & \operatorname{Exan}_{\Lambda}(B,L)\ar[r]
      & \operatorname{Exan}_{\Lambda}(B',L)
  }
\]
où les flèches sont déduites de celles de (18.3.7.1) par fonctorialité.
\end{env}

\begin{proposition}[18.3.8]
\label{0.18.3.8-fr}
Soient $B$ un anneau, $\mathfrak J$ un idéal bilatère de $B$,
$C=B/\mathfrak J$ l'anneau quotient~; $\mathfrak J/\mathfrak J^2$ est alors
canoniquement muni d'une structure de $C$-bimodule. Pour tout $C$-bimodule
$L$, soit $\operatorname{Hom}_C(\mathfrak J/\mathfrak J^2,L)$ le groupe
additif des homomorphismes de $C$-bimodules de $\mathfrak J/\mathfrak J^2$
dans $L$. On définit alors un isomorphisme canonique de groupes commutatifs
\begin{equation}
\label{0.18.3.8.1-fr}
\tag{18.3.8.1}
  \eta_L:\operatorname{Hom}_C(\mathfrak J/\mathfrak J^2,L)
  \xrightarrow{\sim}\operatorname{Exan}_B(C,L)
\end{equation}
en faisant correspondre à tout $C$-homomorphisme
$w:\mathfrak J/\mathfrak J^2\to L$ (qui est \emph{a fortiori} un
$B$-homomorphisme) la classe de l'extension
$(B/\mathfrak J^2)\oplus_{(\mathfrak J/\mathfrak J^2)}L$ déduite de
l'extension $B/\mathfrak J^2$ de $C$ par $\mathfrak J/\mathfrak J^2$ au moyen
de $w$~; l'isomorphisme réciproque fait correspondre à la classe d'une
$B$-extension $E$ de $C$ par $L$ l'homomorphisme $w$ tel que le composé
$\mathfrak J\to\mathfrak J/\mathfrak J^2\xrightarrow{w}L$ soit la restriction
à $\mathfrak J$ de l'homomorphisme structural $B\to E$.
\end{proposition}

\begin{proof}
Soit $F=B/\mathfrak J^2$, qui est une $B$-extension de $C$ par
$\mathfrak J/\mathfrak J^2$. Pour toute $B$-extension
\[
  0\longrightarrow L\xrightarrow{j}E\xrightarrow{p}C\longrightarrow0
\]
de $C$ par $L$, l'homomorphisme structural $f:B\to E$ est tel que le composé
$B\xrightarrow{f}E\xrightarrow{p}C$ soit l'homomorphisme canonique
$B\to B/\mathfrak J=C$. Comme l'image de $\mathfrak J$ par $p\circ f$ est
nulle, $f(\mathfrak J)$ est contenu dans le noyau de $p$, c'est-à-dire
$j(L)$, et comme $j(L)$ est de carré nul, on a $f(\mathfrak J^2)=0$~; donc
$f$ se factorise en
$B\to B/\mathfrak J^2=F\xrightarrow{u}E$, et si $w$ est la restriction de
$u$ à $\mathfrak J/\mathfrak J^2$, on a un diagramme commutatif
\[
  \xymatrix@C=2pc@R=1.5pc{
    0\ar[r] & \mathfrak J/\mathfrak J^2\ar[r]\ar[d]^w
      & F\ar[r]^g\ar[d]^u
      & C\ar[r]\ar[d]^{1_C}\ar@{}[d]|{\wr} & 0\\
    0\ar[r] & L\ar[r]_j & E\ar[r]_p & C\ar[r] & 0
  }
\]

autrement dit $(u,1_C,w)$ est un morphisme d'extensions (18.2.4). On en
déduit un morphisme d'extensions
\[
  \xymatrix@C=2pc@R=1.5pc{
    0\ar[r] & L\ar[r]\ar[d]^{1_L}\ar@{}[d]|{\wr}
      & E'\ar[r]\ar[d]^{u'}
      & C\ar[r]\ar[d]^{1_C}\ar@{}[d]|{\wr} & 0\\
    0\ar[r] & L\ar[r]_j & E\ar[r]_p & C\ar[r] & 0
  }
\]

